\documentclass[11pt]{article}

\author{Math 1030}
\date{Due Friday, April 3 by 5pm} 
\title{Homework 7}

\usepackage{graphicx,xypic}
\usepackage{amsthm}
\usepackage{amsmath,amssymb}
\usepackage{amsfonts}
\usepackage{xcolor}
\usepackage[margin=1in]{geometry}
\usepackage[shortlabels]{enumitem}
\newtheorem{problem}{Problem}
\renewcommand*{\proofname}{{\color{blue}Solution}}


\setlength{\parindent}{0pt}
\setlength{\parskip}{1.25ex}


\begin{document}

\maketitle

% DON'T LEAVE THIS BLANK! 
{\bf\Large Your Name:} 

% DON'T FORGET YOUR COLLABORATORS! 
Collaborator names: 

\vspace{.3in}
Topics covered: graph coloring, planar graphs

Instructions {\bf (read these carefully!)}: 
\begin{itemize}
\item This assignment must be submitted on Gradescope by the due date. 
\item If you collaborate with other students (this is encouraged),  list your collaborators above. 
\item Homework is graded anonymously, so avoid putting your name on each page of the assignment.
\item If you are stuck, please ask for help (from me, Rafi, or a classmate). Use Ed discussions!  
\item You may freely use any fact proved in class (mention it in the appropriate spot). You \emph{may not} freely use facts from the book or other sources.
\item As a general rule, try to restrict your solution to each problem to a single page. Often solutions can be short, while still being complete. If your solution is very long, think more about how you might express it more concisely.
\end{itemize}


\pagebreak 

\begin{problem}
Prove that a planar graph has at least four vertices with degree $\le5$. \footnote{We proved that there is at least one vertex. Observe that the stronger statement cannot be true as stated since it doesn't hold for $K_3$ (!).  Part of your solution should include the correct hypotheses and you should make clear where those hypotheses are needed in the argument.}
\end{problem}

\begin{proof}


\end{proof}

\pagebreak 

\begin{problem}
Find the error in the following incorrect proof of Brooks' Theorem: If $G\neq C_{2k+1}$ and $G\neq K_n$, then $\chi(G)\le\Delta(G)$. A complete solution requires not only pointing out the error, but also explaining why it's wrong (with an example). 

\begin{quotation}
We induct on the number of vertices $n$. The statement holds for $n=1$, trivially. For the induction step, suppose that $G$ is not complete or an odd cycle. Since $\kappa(G)\le\delta(G)$ (where $\kappa$ is the vertex connectivity and $\delta$ is the minimal vertex degree), the graph $G$ has a vertex cut $S$ of size at most $\Delta(G)$. Let $G_1,\ldots,G_m$ be the components of $G\setminus S$, and let $H_i\subset G$ be the subgraph generated by $V(G_i)\cup S$. By the induction hypothesis, each $H_i$ is $\Delta(G)$-colorable. Permute the names of the colors used on these subgraphs to agree on $S$. This yields a coloring of $G$ with $\Delta(G)$ colors. 
\end{quotation} 
\end{problem}

\begin{proof}

\end{proof}

\pagebreak 

\begin{problem}
Given a set of lines in the plane with no three meeting at a point, form a graph $G$ whose vertices are the intersections of the lines, with two vertices adjacent if they appear consecutively on one of the lines. Prove that $\chi(G)\le3$. 
\footnote{Hint: use a greedy coloring with an appropriate vertex ordering. } 
\end{problem}

\begin{proof}

\end{proof}

\pagebreak 

\begin{problem}
Let $G\subset\mathbb R^2$ be a planar graph with a fixed planar embedding. Prove that $G$ is bipartite if and only if each face of $G$ has an even number of sides. 
\end{problem}

\begin{proof}

\end{proof}

\pagebreak 

\begin{problem}Let $G$ be a graph whose complement is bipartite. Prove that the chromatic number of $G$ is equal to the \emph{clique number}, which is the maximum $n$ so that $K_n$ is a subgraph of $G$. \footnote{Hint: look to apply K\"onig's theorem (!)} 
\end{problem}

\begin{proof}

\end{proof}

\pagebreak 

\begin{problem}
Prove the $6$-color theorem by induction (and a fact from class, or even a problem from this assignment). 
\end{problem}

\begin{proof}

\end{proof}

\pagebreak 

\begin{problem}[Bonus]
Look up the Moser spindle (it is a graph). Prove it has chromatic number $4$. Read about the Hadwiger–Nelson problem and explain the relevance of the Moser spindle to this problem. 
\end{problem}

\begin{proof}

\end{proof}





\end{document}